\section{Retained domains and coordinate dependence in P-split}
\label{app:p-split}
Kronqvist, Misener, and Tsay's P-split construction splits convex additive
constraints by groups of variables, introduces epigraph auxiliaries for
the groups, and convexifies the resulting bounded auxiliary disjunction
\citep[Sections 2--4]{KronqvistEtAl2026}. The original domain $x\in X$
remains a global constraint. The bounds are computed over this retained
compact convex $X$, and identical or scaled group functions share an
auxiliary under the paper's minimal-sharing Assumption 3 and Remark 1.
In the following comparisons, ``basic'' means these epigraph links,
exact global interval bounds, and the exact auxiliary disjunction hull.
Extra original--auxiliary links change this model.

\subsection{A domain counterexample and a local nonexactness criterion}
Definition 4 of \citet{KronqvistEtAl2026} calls disjuncts fully disjoint
when they are pairwise disjoint, and their Theorem 6 states universal
nonexactness when the bounds are additive and the constraint functions
strictly convex. With the retained domain, this universal statement
fails. The failure concerns the theorem's scope, not the validity of the
P-split formulations.

\begin{example}[An exact retained-box relaxation]\label{ex:psplit-exact-box}
On $X=[0,3]\times[-1,1]$, use
\[
 D_0=\{(t,w)\in X:t^2+w^2\le1\},\qquad
 D_3=\{(t,w)\in X:(t-3)^2+w^2\le1\}.
\]
The disjuncts are nonempty, disjoint, and defined by strictly convex
quadratics, with additive component bounds on the box. Together they
contain all four vertices of $X$. Hence $\conv(D_0\cup D_3)=X$, and every convex
valid relaxation retaining $X$, including every P-split partition, is exact.
For a direct full-split lift use shared $w^2$ and set
\[
 (\alpha_0,\alpha_3,\alpha_w)=(3t,9-3t,1).
\]
This lies in the auxiliary hull as the combination of $(0,9,1)$ and
$(9,0,1)$ with weights $1-t/3,t/3$. The links follow from
$t^2\le3t$, $(t-3)^2\le9-3t$, and $w^2\le1$. In dimension $n\ge2$
replace $w^2$ by $(n-1)^{-1}\sum_{i=2}^n x_i^2$ on
$[0,3]\times[-1,1]^{n-1}$. The same box-vertex argument applies,
with rational positive definite quadratics and one constraint per disjunct.
\end{example}
The source calls its ``few constraints'' Assumption 4 technically
unnecessary, and the last construction also fits its intended
large-dimension setting. The proof of its Theorem 6 claims strict Jensen
slack in every component. Strict convexity supplies that slack only when
the two arguments differ, and along the top hull segment in the example
$w$ is fixed at one. Moreover, a neighborhood allowed by epigraph slack
need not remain in $X$: an outward step through its top facet is
forbidden.

\begin{proposition}[A sufficient directional repair]\label{prop:psplit-repair}
Let $X$ and $Q$ be convex, $A$ finite, and $F_a$ continuous convex. Define
$R=\{x\in X:\exists\alpha\in Q,F_a(x)\le\alpha_a\ (a\in A)\}$.
Assume $R$ contains the entire true feasible set, and suppose true
feasible points $x^0,x^1$ have lifts $F(x^0),F(x^1)\in Q$.
At $x^\theta=(1-\theta)x^0+\theta x^1$, $0<\theta<1$, let
$\Delta_a=(1-\theta)F_a(x^0)+\theta F_a(x^1)-F_a(x^\theta)$.
If a hull-supporting inequality $c^Tx\le\beta$ is tight at $x^\theta$
and a direction $v$ has $c^Tv>0$, $x^\theta+\epsilon v\in X$ for all
small $\epsilon\ge0$, and
$F_a(x^\theta+\epsilon v)\le F_a(x^\theta)$ for every zero-slack link
and all small $\epsilon\ge0$, then $R$ strictly contains the true hull.
\end{proposition}
\begin{proof}
Keep $\alpha=(1-\theta)F(x^0)+\theta F(x^1)\in Q$. By continuity,
every positive slack persists for a sufficiently small step, and since
there are finitely many links, one common positive radius suffices. The
assumptions preserve the remaining links and the domain, while the
support violation is $\epsilon c^Tv>0$. If varying links have directional Lipschitz constants
$L_a$ on a common radius, imposing $\epsilon L_a\le\Delta_a$ for every
positive-slack link gives an explicit allowable step within that radius.
\end{proof}

\subsection{Exact basic relaxation of two translated balls}
\begin{proposition}[Projection and Hausdorff error]\label{prop:psplit-balls}
For $n\ge2$, radius $r>0$ and separation $d>2r$, let $x=(t,w)$ and
\[
 D_0=\{t^2+\|w\|_2^2\le r^2\},\qquad
 D_d=\{(t-d)^2+\|w\|_2^2\le r^2\},
 \quad X=[-r,d+r]\times[-r,r]^{n-1}.
\]
Their hull is $C=[0,d]e_1+r\mathbb B_2^n$. The basic full coordinate
split, sharing the transverse square auxiliaries, projects exactly to
\begin{equation}\label{eq:psplit-ball-projection}
 R=\{(t,w):\|w\|_2\le r,\quad
                2(t-d/2)^2+\|w\|_2^2\le2(d/2+r)^2\}.
\end{equation}
The two-group partition $\{t\},\{w_1,\ldots,w_{n-1}\}$ gives the same
set. With $e_*=[(d/2+r)^2-r^2/2]^{1/2}-d/2>0$,
\begin{equation}\label{eq:psplit-ball-distance}
 d_H(R,C)=\sqrt{r^2+e_*^2}-r,
 \qquad d_H(R,C)/r\longrightarrow\sqrt2-1\quad(d/r\to\infty).
\end{equation}
\end{proposition}
\begin{proof}
Put $U=(d+r)^2$. The auxiliaries are $a\ge t^2$, $b\ge(t-d)^2$,
$c_i\ge w_i^2$, with $0\le a,b\le U$, $0\le c_i\le r^2$.
The auxiliary disjunction is $a+\sum_i c_i\le r^2$ or
$b+\sum_i c_i\le r^2$. Its hull is exactly
\begin{equation}\label{eq:psplit-aux-hull}
 0\le a,b\le U,\quad c_i\ge0,\quad\sum_i c_i\le r^2,
 \qquad a+b+\sum_i c_i\le U+r^2.
\end{equation}
These inequalities hold in each disjunct. For sufficiency fix $c$ and put
$h=r^2-\sum_i c_i\in[0,r^2]$. At this fixed $c$, the union in $(a,b)$
is $[0,h]\times[0,U]\cup[0,U]\times[0,h]$. Its hull is the square
cut by $a+b\le U+h$: every vertex of that truncated square lies in
the union. This proves sufficiency with the same $c$ in each constituent.
The set \eqref{eq:psplit-aux-hull} is downward closed, so its epigraph
projection is obtained by substituting the actual squares. Completing the
square gives \eqref{eq:psplit-ball-projection}; its cylinder and ellipsoid
already imply all the retained box bounds. In the two-group partition
replace $\sum_i c_i$ by $c\ge\|w\|^2$. Its global upper bound
$(n-1)r^2$ is reduced to $r^2$ by either disjunct, and the same proof
applies.
The additive-bound refinement property
\citep[Theorem 4 and Corollary 3]{KronqvistEtAl2026} puts this full-split
set inside every basic coarser coordinate partition for this box.

Between the centers every point of the cylinder is in $C$. By symmetry,
consider only $t\le0$ and put $\rho=\|w\|\in[0,r]$. The leftmost allowed
coordinate is $t=-e(\rho)$ with
$e(\rho)=\sqrt{(d/2+r)^2-\rho^2/2}-d/2>0$.
The squared distance to the nearer center is
\[
 e(\rho)^2+\rho^2=(d/2+r)^2+d^2/4+\rho^2/2
                         -d\sqrt{(d/2+r)^2-\rho^2/2}.
\]
Its derivative with respect to $\rho^2$ is
$1/2+d/[4\sqrt{(d/2+r)^2-\rho^2/2}]>0$. The maximum occurs at
$\rho=r$, where the distance outside the capsule is
$\sqrt{r^2+e_*^2}-r$. At $\rho=0$ that distance is zero, so taking
positive parts loses nothing. Since $C\subset R$, the directed and
ordinary Hausdorff distances coincide. Finally $e_*/r\to1$, proving the limit.
\end{proof}

\subsection{An unbounded coordinate effect with rational data}
First take unit balls centered at $0$ and $(D,D)$, $D\ge5$, retaining
$[-1,D+1]^2$. The basic split introduces four separate squares
$a_i\ge x_i^2$, $b_i\ge(x_i-D)^2$ with upper bound $U=(D+1)^2$,
and convexifies $a_1+a_2\le1$ or $b_1+b_2\le1$.
At $p=(2D/3,D/3)$ the true square vectors $a,b$ have
$2a_i,2b_i\le8D^2/9\le U$. Their average representation
$(a,b)=\tfrac12(0,2b)+\tfrac12(2a,0)$ proves feasibility in the
auxiliary hull. Projection onto the center segment is $(D/2,D/2)$,
so the Hausdorff error is at least $D/(3\sqrt2)-1$.
Under the orthogonal change $t=(x_1+x_2)/\sqrt2$,
$w=(x_1-x_2)/\sqrt2$, the retained domain is exactly
\[
 -\sqrt2\le t+w,t-w\le\sqrt2(D+1),
\]
a diamond, and the centers are $0$ and $\sqrt2D e_t$. The shared $w^2$
auxiliary is at most one throughout either disjunct, hence $|w|\le1$
in the relaxation. For $t\le0$, the domain implies
$-t+|w|\le\sqrt2$, giving
$t^2+w^2\le(\sqrt2-|w|)^2+w^2\le2$; the convex quadratic in
$|w|\in[0,1]$ has its maximum at an endpoint. The right end is
identical, and between the centers the cylinder lies in the capsule.
The new error is therefore at most $\sqrt2-1$. No additive-bound
hierarchy is asserted on the transformed diamond.

\begin{proposition}[Rational coordinate comparison]\label{prop:psplit-rational}
For rational $D\ge2$, unit balls with centers $0,v=(3D,4D)$ in
$X=[-1,3D+1]\times[-1,4D+1]$ have basic original-coordinate Hausdorff
error at least $4D/5-1$. Under the fixed rational orthogonal map
\begin{equation}\label{eq:psplit-rational-map}
 t=(3x_1+4x_2)/5,\qquad w=(4x_1-3x_2)/5,
\end{equation}
using the exact transformed domain and tight recomputed bounds, the error
is at most $\sqrt2-1$.
\end{proposition}
\begin{proof}
The witness $p=(2D,4D/3)$ has each coordinate at one or two thirds of
$v_i$. Twice each original square is at most $8v_i^2/9\le(v_i+1)^2$,
so the preceding half-weight lift is feasible. Its transformed coordinates
are $t(p)=34D/15\in[0,5D]$ and $w(p)=4D/5$, proving the lower bound.
In aligned coordinates the common transverse auxiliary again forces
$|w|\le1$. For $t\le0$, the inverse map gives $x_1\le4/5$ and
$x_2\le3/5$, while the retained domain gives $x_1,x_2\ge-1$.
Thus $t^2+w^2=x_1^2+x_2^2\le2$. For $t\ge5D$, apply the same argument
to $x-v$: its two components lie in $[-4/5,1]$ and $[-3/5,1]$.
The error outside either end is at most $\sqrt2-1$, and between the
centers it is zero. The map has determinant $-1$; changing the sign of $w$
gives a proper rotation with identical conclusions. All bounds and domain
inequalities remain rational for rational $D$.
\end{proof}

\subsection{The exact auxiliary image is still different from the graph}
\begin{proposition}[Strongest auxiliary-only comparison]
\label{prop:psplit-exact-image}
In either original-coordinate example, let
$F(x)=(x_i^2,(x_i-v_i)^2)_{i=1}^2$. Every convex auxiliary set $Q$
containing $F(0)$ and $F(v)$ preserves the original witness in
\[
 R(Q)=\{x\in X:\exists\alpha\in Q,F(x)\le\alpha\}.
\]
In particular the lower error survives $Q=\conv F(D_0\cup D_v)$.
In aligned coordinates, however, the exact feasible auxiliary image hull
for $F'(t,w)=(t^2,(t-d)^2,w^2)$ gives exactly the true capsule, for $d>2$.
This holds with either exactly transformed retained domain above.
\end{proposition}
\begin{proof}
Convexity forces $[F(0)+F(v)]/2\in Q$, whose entries are $v_i^2/2$.
At the witness, both coordinate squares are at most $4v_i^2/9$, so they
satisfy all epigraph links to that point. The argument uses only the two
feasible center images, so every valid convex auxiliary-only
strengthening of the same epigraph map retains the witness.

For alignment write $(a,b,c)=F'(t,w)$ and set
\[
 f(c)=d^2+1-c+2d\sqrt{1-c},\qquad0\le c\le1.
\]
At an actual feasible point, $t\in[-q,q]\cup[d-q,d+q]$ with
$q=\sqrt{1-c}$, so $a,b\le(d+q)^2=f(c)$. The function $f$ is
concave and decreasing. Its two hypographs therefore contain the whole
convex auxiliary image hull. If the epigraph links are feasible there,
then $|w|\le1$ and, by monotonicity,
\[
 |t|\le d+\sqrt{1-w^2},\qquad |t-d|\le d+\sqrt{1-w^2}.
\]
Their intersection is precisely
$-\sqrt{1-w^2}\le t\le d+\sqrt{1-w^2}$, the capsule.
Conversely the epigraph relaxation is convex and contains both balls,
so it contains their hull. The retained domain contains that hull,
which proves equality. In fact, $0\le c\le1$ and $a,b\le f(c)$ alone
suffice: each upper inequality has a conic representation with
$s\ge0$, $s^2+c\le1$ and, respectively,
$a+c-d^2-1\le2ds$ or $b+c-d^2-1\le2ds$.
\end{proof}
The exact auxiliary image hull retains only values of the chosen functions;
it is not the hull of their full feasible graph in $(x,\alpha)$.
Original-space cuts, constraints coupling original and auxiliary variables,
or identities replacing the epigraph links can remove the witness.
These calculations establish explicit coordinate effects in this family.
They give no general algorithm for choosing good coordinates and no
assertion about every strengthened P-split formulation.
